
\subsection{Holographic Transformation}
To introduce the idea of holographic transformation,
it is convenient to consider bipartite graphs.
For a general graph,
we can always transform it into a bipartite graph while preserving the Holant value,
as follows.
For each edge in the graph,
we replace it by a path of length two.
(This operation is called the \emph{2-stretch} of the graph and yields the edge-vertex incidence graph.)
Each new vertex is assigned the binary \textsc{Equality} signature $=_2$. Thus, we have $\holant{=_2}{\mathcal{F}}\equiv_T \Holant(\mathcal{F})$.


For an invertible $2$-by-$2$ matrix $T \in {\bf GL}_2({\mathbb{C}})$
 and a signature $f$ of arity $n$, written as
a column vector (covariant tensor) $f \in \mathbb{C}^{2^n}$, we denote by
$Tf = T^{\otimes n} f$ the transformed signature.
  For a signature set $\mathcal{F}$,
define $T\mathcal{F} = \{T f \mid  f \in \mathcal{F}\}$ the set of
transformed signatures.
For signatures written as
 row vectors (contravariant tensors) we define
$f T^{-1}$ and  $\mathcal{F} T^{-1}$ similarly.
Whenever we write $T f$ or $T \mathcal{F}$,
we view the signatures as column vectors;
similarly for $f T^{-1}$ or $\mathcal{F} T^{-1}$ as row vectors.
We can also represent $Tf$ as the matrix $M_{S_k}(Tf)$ with the assignments of variables in $S_k$ as row index and the assignments of the other $n-k$ variables as column index. 
Then,  we  have $M_{S_k}(Tf)=T^{\otimes k}M_{S_k}(f)(T^{\tt T})^{\otimes n-k}$. Similarly, $M_{S_k}(fT^{-1})=({T^{-1}}^{\tt T})^{\otimes k}M_{S_k}(f)(T^{-1})^{\otimes n-k}$.
%In the special case of the Hadamard matrix
%$H_2 = \frac{1}{\sqrt{2}} \left[\begin{smallmatrix} 1 & 1 \\ 1 & -1 \end{smallmatrix}\right]$,
%we also define $\widehat{\mathcal{F}} = H_2  \mathcal{F}$.
%Note that $H_2$ is orthogonal.
%Since constant factors are immaterial, for convenience we sometime
%drop the factor $\frac{1}{\sqrt{2}}$ when using $H_2$. 

Let $T \in {\bf GL}_2({\mathbb{C}})$.
The holographic transformation defined by $T$ is the following operation:
given a signature grid $\Omega = (H, \pi)$ of $\holant{\mathcal{F}}{\mathcal{G}}$,
for the same bipartite graph $H$,
we get a new signature grid $\Omega' = (H, \pi')$ of $\holant{\mathcal{F} T^{-1}}{T \mathcal{G}}$ by replacing each signature in
$\mathcal{F}$ or $\mathcal{G}$ with the corresponding signature in $\mathcal{F} T^{-1}$ or $T \mathcal{G}$.

\begin{theorem}[\cite{Valiant08}]\label{thm: holographic transformation}
 For every $T \in {\bf GL}_2({\mathbb{C}})$,
  $\Holant(\mathcal{F} \mid \mathcal{G}) \equiv_T \Holant(\mathcal{F} T^{-1} \mid T \mathcal{G}).$
\end{theorem}

Therefore,
a holographic transformation does not change the complexity of the Holant problem in the bipartite setting. 
A particular useful holographic transformation is the $K$-transformation.
Direct calculation shows $(=_2)K=(\neq_2)$, so
\(
\holant{=_2}{\mathcal F}
\equiv_T
\holant{\neq_2}{K^{-1}\mathcal F}.
\)
Accordingly, throughout this paper
\[
\widehat f:=K^{-1}f,
\qquad
\widehat{\mathcal F}:=K^{-1}\mathcal F.
\]
We also call
$\holant{\neq_2}{\widehat{\mathcal F}}$ the
{$K$-Holant} formulation of $\Holant(\mathcal F)$.


\begin{lemma}\label{lm: conjugation and K-trans}
Let $f$ be a complex-valued signature of arity $n$.  For every
$\alpha\in\{0,1\}^n$,
\(
\widehat{\overline f}(\alpha)
=\overline{\widehat f(\overline\alpha)}.
\)
\end{lemma}
\begin{proof}
Using $K^{\mathtt T}=N_2K^{-1}$, we have
\(
\widehat{\overline f}
=(K^{-1})^{\otimes n}\overline f
=\overline{(K^{\mathtt T})^{\otimes n}f}
=N_2^{\otimes n}\overline{\widehat f},
\)
which is the stated entrywise identity.
\end{proof}

The preceding lemma immediately gives
the following equivalence.
\begin{lemma}\label{lm: conjugate closed and ars closed}
$\mathcal F$ is conjugate closed if and only if
$\widehat{\mathcal F}$ is arrow-reversal closed.
\end{lemma}

